\section{Related Work}
\noindent\textbf{State representation for control.}
In learning end-to-end visuomotor policies, the policy is typically parameterized as a neural network that takes image observation as the input representation~\cite{reed2022generalist,brohan2022rt,walke2023bridgedata}. Due to the lack of inductive bias, these approaches require training on a large number of demonstration trajectories, which is expensive to collect. On the other hand, different structured representations are proposed to improve the data efficiency, such as key points~\cite{qin2020keto,manuelli2020keypoints}, mesh~\cite{lin2021VCD,huang2022medor}, or neural 3D representation~\cite{li20213d,rashid2023language}. However, prior structures often limit the policy to specific tasks. In contrast, we propose to utilize future trajectories of arbitrary points in the image as additional input to the policy, making minimal assumptions about the task and environment. We demonstrate its wide application to a set of over 130 language-conditioned manipulation tasks. In navigation and locomotion, it is common to construct policies that are guided by the future trajectories of the robot~\cite{aguiar2007trajectory,peng2018deepmimic}. In manipulation, some works have explored flow-based guidance~\cite{goyal2022ifor,seita2023toolflownet,gu2023rt}. However, prior works only track task-specific points, such as the end-effector of the robot, or the human hand. Instead, our approach works with arbitrary points, including points on the objects, thus providing richer information in the more general settings. Finally, \citet{vecerik2023robotap} proposes to utilize any-point tracking for few-shot policy learning. This approach does not learn trajectory models from data, but instead mainly uses the tracker to perform visual servoing during test time. This design choice requires more instrumentation, such as separating the task into multiple stages, predicting the goal locations of the points for each stage, and running the tracker during inference time. In contrast, we present a much simpler framework, enabling application in more diverse settings.

\noindent\textbf{Video pre-training for control.} Videos contain rich information about behaviors and dynamics, which can help policy learning. However, video pre-training remains challenging due to the lack of action labels. One line of works first learns an inverse dynamics model that predicts the action from two adjacent frames and then labels the videos with pseudo actions~\cite{schmeckpeper2019learning, baker2022vpt, torabi2018behavioral,schmidt2024learning}. However, the inverse dynamics model is often trained on a limited action-labeled dataset and does not generalize well, especially for continuous actions. Prior works have also explored pre-train a feature representation using various self-supervised objectives~\cite{sermanet2017tcn,nair2023r3m,ma2023vip}, but the representation alone does not retain the temporal information in videos. More recently, learning video prediction as pre-training has shown promising results~\cite{du2023unipi,Ko2023Learning,bharadhwaj2023zero,yang2023learning,Escontrela23arXiv_VIPER}. During policy learning, a video prediction model is used to generate future subgoals and then a goal-conditioned policy can be learned to reach the sub-goal. However, video prediction models often result in hallucinations and unrealistic physical motions. Furthermore, video models require extensive computation, which is an issue, especially during inference time. In contrast, our method models the trajectories of the points, naturally incorporating inductive bias like object permanence while requiring much less computation. This enables our trajectory models to be run closed-loop during policy execution. Furthermore, the trajectories provide dense guidance to the policy as a motion prior.

\noindent\textbf{Learning from Human Videos.} Of particular interest, numerous prior works have proposed learning from the rich source of human videos~\cite{xiong2021learning,shao2021concept2robot, mendonca2023structured,bahl418affordances,shaw2023videodex,wang2023mimicplay}. However, these works often explicitly extract the hand pose or contact regions from the human videos, thereby losing information about the dynamics of the remaining objects. In contrast, our method models the trajectories of arbitrary points and can learn from both human videos and videos of a different robot.